% Irodov Part V §5.3 originals 131--142. Original PDF printed 209--211,
% source page 210 visually checked: #139 prints 10.0 m, #140 prints 670.8 mm,
% both inconsistent with answer; variants explicitly preserved.
\diffractionproblem{131}{209}
图5.26：平面单色光 $\lambda$ 垂直入射玻璃相位光栅，
材料折射率 $n$，相邻沟槽与凸脊等宽均为 $a$、刻槽深度 $h$。
求：(a) 使零级夫琅禾费极大强度为零的深度；
(b) 此时一级衍射角。
\exsolution
沟槽与未刻区的光程差
$\Delta=(n-1)h$，相位差
$\delta=2\pi(n-1)h/\lambda$。
零级复振幅与
$\tfrac12(1+e^{i\delta})$ 成正比，
因此零级为零要求
$\delta=(2j+1)\pi$，即
\[
\boxed{h_j=\frac{(2j+1)\lambda}{2(n-1)}},\qquad j=0,1,\ldots.
\]
光栅周期 $d=2a$，一级主极大
$d\sin\theta_1=\lambda$，所以
\[
\boxed{a\sin\theta_1=\lambda/2}.
\]
\excheck 相位光栅的零级消光来自等面积两区域\emph{复振幅}相反，不是吸收光损失。

\diffractionproblem{132}{209}
图5.27：波长
$\lambda=0.55\,\mathrm{\mu m}$ 的平面光穿过装有水的超声槽，
槽内驻波频率
$\nu=4.7$MHz，后接焦距
$f=35$cm 的物镜。
焦面相邻主极大距离
$\Delta x=0.60$mm，求水中超声波传播速度。
\exsolution
超声形成沿传播方向周期性折射率调制，其有效光栅周期与声波长相同：
$d=v_s/\nu$。
近轴一级衍射角差
$\Delta\theta\simeq\lambda/d$，
焦面间距
$\Delta x\simeq f\Delta\theta$。
故
\[
\boxed{v_s=\frac{\lambda\nu f}{\Delta x}}
=\frac{(0.55\times10^{-6})(4.7\times10^6)(0.35)}
{0.60\times10^{-3}}
\simeq1.5\,\mathrm{km/s}.
\]
\excheck 源题说明是超声驻波；本模型的有效周期按原书声波空间周期处理，不能凭“驻波”一词无说明地再除以2。

\diffractionproblem{133}{209}
用迈克耳孙方法测量双星角距 $\psi$：
在望远镜物镜前放间距可调的两条平行窄缝，
缩小缝距至
$d=95$cm 时首次观察到条纹可见度变差。
取波长 $\lambda=0.55\,\mathrm{\mu m}$，求 $\psi$。
\exsolution
双星两点源各产生一套双缝条纹；
它们对应的中心相位差为
$\Delta\phi=(2\pi d/\lambda)\psi$
（小角近似）。
首次最差可见度发生在两套条纹互移半周期：
$\Delta\phi=\pi$，
即
\[
\boxed{\psi=\frac{\lambda}{2d}}
=\frac{0.55\times10^{-6}}{2(0.95)}
\simeq2.9\times10^{-7}\,\mathrm{rad}.
\]
用 $1$rad $=206265''$，
\[
\boxed{\psi\simeq0.060''}.
\]
\excheck 两套条纹来自\emph{不相干}恒星，
相互不能相干叠加振幅，但两套强度花样错位会降低总可见度。

\diffractionproblem{134}{210}
透射光栅周期
$d=1.50\,\mathrm{\mu m}$，
波长 $\lambda=530$nm。
求最高可见级主极大对应角散度
$D=\mathrm d\theta/\mathrm d\lambda$，
单位为角分/nm：
(a) 垂直入射；
(b) 入射角 $\theta_0=45^\circ$。
\exsolution
有向光栅方程
$d(\sin\theta-\sin\theta_0)=m\lambda$，
对波长微分：
\[
\boxed{|D|=\frac{|m|}{d|\cos\theta|}}.
\]
(a) 正入射最高级
$|m|=\lfloor d/\lambda\rfloor=2$，
$\sin\theta=2\lambda/d$，
所以
\[
\boxed{D_a=
\frac2{d\sqrt{1-(2\lambda/d)^2}}
\simeq6.5\,\text{角分/nm}}.
\]
(b) 因角度向另一侧延伸，最高绝对级是
$m=-4$，
$\sin\theta=\sin45^\circ-4\lambda/d$，
\[
\boxed{D_b=
\frac4{d\sqrt{1-(\sin45^\circ-4\lambda/d)^2}}
\simeq13\,\text{角分/nm}}.
\]
\excheck 将 rad/m 换成角分/nm：
乘 $3437.75\times10^{-9}$；不应仅报告无量纲13。

\diffractionproblem{135}{210}
波长为 $\lambda$ 的平面光垂直入射周期 $d$ 的光栅。
求主极大角散度 $D=\mathrm d\theta/\mathrm d\lambda$
用衍射角 $\theta$ 表示的公式。
\exsolution
第 $m$ 级光栅方程
$d\sin\theta=m\lambda$。
对 $\lambda$ 微分
$d\cos\theta\,\mathrm d\theta=m\,\mathrm d\lambda$，
从而
$D=m/(d\cos\theta)$。
再由 $m/d=\sin\theta/\lambda$，
得到
\[
\boxed{D=\frac{\tan\theta}{\lambda}}.
\]
\excheck 固定的是\emph{衍射级次} $m$；若改变级次而非波长，上式不是同一个偏导数。

\diffractionproblem{136}{210}
波长
$\lambda=589.0$nm 的光垂直入射光栅，
周期 $d=2.5\,\mathrm{\mu m}$，
总刻线数 $N=10000$。
求二级主极大的全角宽（两侧最近零点之间）。
\exsolution
二级方向满足
$\sin\theta_2=2\lambda/d$。
$N$ 缝干涉因子的最近零点在
$d\sin\theta=(2\pm1/N)\lambda$，
故峰两侧的全角宽
\[
\boxed{\Delta\theta\simeq
\frac{2\lambda}{Nd\cos\theta_2}
=\frac{2\lambda}
{Nd\sqrt{1-(2\lambda/d)^2}}}.
\]
代数值
$\Delta\theta\simeq5.3\times10^{-5}$rad，
\[
\boxed{\Delta\theta\simeq11''}.
\]
\excheck 题中“angular width”采用全主峰零点间宽度；若定义半宽则差一个因子2。

\diffractionproblem{137}{210}
证明光垂直入射宽度为 $\ell$ 的透射光栅时，
其最大分辨本领不能超过 $\ell/\lambda$。
\exsolution
波长分辨本领由瑞利判据给出
$\mathcal R=\lambda/\delta\lambda=mN$，
其中 $N$ 是照亮的刻线数、$m$ 为衍射级次。
若光栅周期为 $d$，则 $N\simeq\ell/d$。
另一方面衍射必须满足
$|m\lambda/d|\le1$，
即 $|m|\le d/\lambda$。
两者结合
\[
\boxed{\mathcal R_{\max}=mN
\le\frac d\lambda\frac\ell d
=\frac\ell\lambda}.
\]
\excheck 极限只是几何与衍射级次的上界，实际仪器还受光栅效率、缺陷、像差影响。

\diffractionproblem{138}{210}
用衍射光栅说明：
按瑞利判据刚能分辨的两条相邻谱线的频差
$\delta\nu$ 等于参与干涉的最远两振动传播时间差 $\Delta t$ 的倒数。
\exsolution
对第 $m$ 级，
在最远两光栅缝之间的路程差约
$\Delta\ell=N d\sin\theta
=N m\lambda$，
传播时延
$\Delta t=\Delta\ell/c
=Nm/\nu$。
瑞利判据的相邻谱线刚可分辨时，
$\delta\nu/\nu\simeq1/(mN)$，
故
\[
\boxed{\delta\nu=
\frac{\nu}{mN}
=\frac1{\Delta t}}.
\]
\excheck 对有限光栅应取 $N-1$ 个缝间隔，更精确的系数含 $N-1$；
书中公式为刻线多时 $N\gg1$ 的近似。

\diffractionproblem{139}{210}
两条相近的单色谱线波长分别为
$600.000\,\mathrm{nm}$ 与 $600.050\,\mathrm{nm}$，
垂直入射透射光栅。问在怎样的衍射角
$\theta$ 下恰好达到瑞利分辨判据？
\textbf{宽度版本差异：}多个可查的 Mir 等英文题面记光栅总宽
$\ell=10.0\,\mathrm{mm}$；
本项目此前记录的上传扫描底本题面记为
$\ell=10.0\,\mathrm m$。
两者相差1000倍，必须分别计算。
\exsolution
瑞利判据为 $\lambda/\delta\lambda=mN$。
另一方面正入射时 $d\sin\theta=m\lambda$，
刻线数 $N=\ell/d$。
将级次与刻线间距消去，得到
\[
 \boxed{\sin\theta
 =\frac{\lambda^2}{\ell\,\delta\lambda}},
 \qquad
 \lambda\simeq600.0\,\mathrm{nm},\quad
 \delta\lambda=0.050\,\mathrm{nm}.
\]
\textbf{标准英文版本 $\ell=10.0\,\mathrm{mm}$：}
\[
 \sin\theta=\frac{(600.0\times10^{-9})^2}
 {(10.0\times10^{-3})(0.050\times10^{-9})}
 =0.720,
 \qquad
 \boxed{\theta\simeq46.05^\circ\approx46^\circ}.
\]
\textbf{若严格按已记录扫描题面
$\ell=10.0\,\mathrm m$：}
\[
 \sin\theta=7.20\times10^{-4},\qquad
 \boxed{\theta\simeq0.0413^\circ}.
\]
\excheck
维数检查：$\lambda^2/(\ell\,\delta\lambda)$ 无量纲。
按标准英文版 $\ell=10.0\,\mathrm{mm}$ 可复现书末
$46^\circ$；按扫描底本记录 $10.0\,\mathrm m$ 不能。
公开对照可见
\texttt{studiestoday.com/ie-irodov-solutions-chapter-29-diffraction-of-light-part-2}。
这一修订只明确排定两种版本的主次，
不声称在本次重看了用户扫描页。

\diffractionproblem{140}{210}
正入射的透射光栅宽 $\ell=6.5\,\mathrm{cm}$，
每毫米刻有200条线。研究波长
$\lambda=670.8\,\mathrm{nm}$、
双线差 $\delta\lambda=0.015\,\mathrm{nm}$ 的谱线。
求：(a) 第几级起可分辨双线；
(b) 在波长约 $670\,\mathrm{nm}$ 处，
该光栅的最小理论可分辨波长差。
\textbf{单位版本差异：}本项目先前记录的扫描底本
第一处记成“$670.8\,\mathrm{mm}$”，
但其后文又明确要求在 $670\,\mathrm{nm}$ 附近讨论；
另有英文原题转录写作 $670.8\,\mathrm{nm}$。
\exsolution
刻线总数
\[
 N=\ell\frac{200}{\mathrm{mm}}
 =(65\,\mathrm{mm})(200/\mathrm{mm})
 =13000,
 \quad
 d=\frac{1}{200}\,\mathrm{mm}=5.0\,\mathrm{\mu m}.
\]
在第 $m$ 级满足 $\mathcal R=mN$，由
$mN\ge\lambda/\delta\lambda$ 得
\[
 m\ge\frac{670.8}{0.015(13000)}
 \simeq3.440,
 \quad\boxed{m_{\min}=4}.
\]
正入射真实传播级次满足
$m\lambda/d\le1$，
故 $m_{\max}=\lfloor d/\lambda\rfloor=7$。
最小可分辨谱线差
\[
 \boxed{\delta\lambda_{\min}
 =\frac{\lambda}{m_{\max}N}
 \simeq\frac{670\,\mathrm{nm}}{7(13000)}
 =0.00736\,\mathrm{nm}
 \simeq7.4\,\mathrm{pm}}.
\]
\excheck
若按“$670.8\,\mathrm{mm}$”字面处理，
单线波长大于 $5.0\,\mathrm{\mu m}$ 的光栅周期约
$1.34\times10^5$ 倍，不可能形成第一级光谱，
更不可能得到第四级；故记录为扫描版本单位冲突，
不能把毫米和纳米分支的条件混入同一推导。

\diffractionproblem{141}{210}
一片宽
$\ell=10$mm 的透射光栅，在垂直入射时，
钠黄双线 $589.0$nm 与 $589.6$nm 从第五级开始可被分辨。
估计：(a) 光栅周期；
(b) 若仍用同样周期，欲在第三级分辨波长
$\lambda=460.0$nm、间隔
$\delta\lambda=0.13$nm 的双线，光栅宽度至少应取多少？
\exsolution
(a) 临界分辨本领
$mN\simeq\lambda/\delta\lambda$，
取 $m=5$，
$N\simeq589/[5(0.6)]\simeq196$。
周期
\[
\boxed{d\simeq\frac{\ell}{N}
\simeq\frac{10\,\mathrm{mm}}{196}
\simeq0.050\,\mathrm{mm}}.
\]
(b) 第三级
$N_{\rm req}\simeq460/[3(0.13)]
\simeq1180$，
所以
\[
\boxed{\ell_{\min}\simeq N_{\rm req}d
\simeq1180(0.050\,\mathrm{mm})
\simeq6.0\,\mathrm{cm}}.
\]
\excheck 第一级别能分辨是在$m=5$，
并不说明其他更高衍射级的主峰强度足够强；
此处只比较理论分辨本领。

\diffractionproblem{142}{210--211}
石英光谱仪透射光栅宽
$\ell=25$mm、每毫米250条线，
物镜焦距
$f=80$cm。
波长为 $310.154$ 与
$310.184$nm 的双线垂直入射。
求：(a) 在一级、二级谱中像面上的分离距离；
(b) 两级是否可分辨。
\exsolution
光栅周期
$d=1/250\,\mathrm{mm}=4\,\mathrm{\mu m}$，
总刻线数
$N=25(250)=6250$。
双线差
$\delta\lambda=0.030$nm，
第 $m$ 级满足
$d\sin\theta_m=m\lambda$。
焦平面坐标 $x_m=f\tan\theta_m$，
故
\[
\delta x_m\simeq
f\frac{m\,\delta\lambda}{d\cos^3\theta_m}.
\]
代 $m=1,2$：
\[
\boxed{\delta x_1\simeq6\,\mathrm{\mu m}},
\qquad
\boxed{\delta x_2\simeq12\,\mathrm{\mu m}}.
\]
瑞利分辨条件
$\delta\lambda\ge\lambda/(mN)$：
一级阈值约 $0.050$nm（大于0.030nm），不可分辨；
二级约 $0.025$nm（小于0.030nm），可以分辨。
\[
\boxed{\text{一级不能分辨；二级可以分辨。}}
\]
\excheck 判断能否分辨应使用相同的光栅线数与瑞利准则，而不是只看像面双线分离几微米。
